\section{Công trình liên quan}

Trong kỷ nguyên học sâu \cite{5, 15, 16, 31, 50, 52, 53, 54, 55, 56, 57, 62, 63, 65, 67, 68, 69, 70, 71}, mục tiêu cốt lõi trong ZSL \cite{7, 8, 9, 11, 14, 22, 23, 24, 25, 28, 42, 45, 58, 60} là tổng quát hóa kiến thức thị giác-ngữ nghĩa thu được từ các lớp đã thấy sang các lớp chưa thấy. Sự tổng quát hóa này dựa vào các bộ mô tả ngữ nghĩa, thường được biểu diễn dưới dạng thuộc tính \cite{21}, tài liệu \cite{20}, hoặc nhúng từ \cite{36}. Các phương pháp ZSL chính thống có thể được phân loại chung thành các phương pháp sinh (generative methods) và các phương pháp dựa trên nhúng (embedding-based methods). Các phương pháp ZSL sinh thường tổng hợp các đặc trưng thị giác chưa thấy từ các bộ mô tả ngữ nghĩa bằng cách sử dụng các mô hình sinh, cho phép các quy trình phân loại tiêu chuẩn kết hợp các mẫu huấn luyện tổng hợp. Các phương pháp này tận dụng các kiến trúc sinh khác nhau, bao gồm GANs \cite{45, 60, 73}, VAEs \cite{43, 51} và Normalizing Flows \cite{10, 12, 44}. Các phương pháp ZSL dựa trên nhúng nhằm mục đích học một ánh xạ trực tiếp từ không gian thị giác sang không gian ngữ nghĩa, liên kết các biểu diễn hình ảnh với các bộ mô tả lớp. Các phương pháp này thường tận dụng cơ chế chú ý để khoanh vùng các bộ mô tả ngữ nghĩa. Zhu và cộng sự \cite{74} đã giới thiệu một mạng lưới cục bộ hóa đa chú ý có hướng dẫn ngữ nghĩa để trích xuất các đặc trưng chi tiết cục bộ của các bản vá phân biệt. Xu và cộng sự \cite{61} đã đề xuất một mạng lưới nguyên mẫu thuộc tính, trong đó mỗi nguyên mẫu đại diện cho một khái niệm ngữ nghĩa cụ thể và được sử dụng làm bộ lọc 1x1 để tích chập các bản đồ đặc trưng từ ResNet101 nhằm tập trung vào vùng quan tâm. Công trình tiếp theo \cite{33} đã đề xuất sử dụng mô hình GloVe \cite{39} để trích xuất các véc-tơ ngữ nghĩa từ tên thuộc tính (ví dụ: "ngực trắng"), được sử dụng thêm để đại diện cho các nguyên mẫu thuộc tính. Gần đây hơn, Chen và cộng sự \cite{3} đã phát triển một mạng lưới chưng cất ngữ nghĩa tương hỗ để tinh chỉnh sự liên kết giữa các đặc trưng thị giác và ngữ nghĩa.

Mặc dù có những tiến bộ trong các phương pháp dựa trên sinh và nhúng, sự sai lệch ngữ nghĩa vẫn là một thách thức cơ bản trong ZSL \cite{1, 4, 6, 9, 27, 30}. Một số công trình cố gắng tách rời các đặc trưng liên quan đến ngữ nghĩa khỏi các đặc trưng không liên quan đến ngữ nghĩa. DLFZRL \cite{46} là công trình đầu tiên giới thiệu ZSL dựa trên sự tách rời, chia các biểu diễn hình ảnh thành các đặc trưng ngữ nghĩa, không ngữ nghĩa, và không phân biệt. RFF \cite{27} đã áp dụng tối thiểu hóa thông tin tương hỗ để loại bỏ các chi tiết thị giác dư thừa, trong khi SDGZSL \cite{9} sử dụng phân tích tương quan tổng thể để tách các thành phần nhất quán về mặt ngữ nghĩa và không liên quan đến ngữ nghĩa trong không gian đặc trưng. FREE \cite{1} đã tinh chỉnh các đặc trưng thị giác thông qua một mất mát trung tâm biên tự thích ứng để tăng cường sự liên kết ngữ nghĩa. Dis-VAE \cite{30} đã giới thiệu một chiến lược tái tổ hợp lô để chưng cất các đặc trưng cụ thể của danh mục để nhận dạng chính xác hơn. Phương pháp ZSLViT dựa trên ViT gần đây \cite{4} giới thiệu một framework ZSL dựa trên vision transformer giúp hợp nhất các token không liên quan đến ngữ nghĩa để cải thiện sự liên kết thị giác-ngữ nghĩa trong việc học các biểu diễn ZSL. Mặc dù có những nỗ lực trong việc giảm thiểu vấn đề sai lệch ngữ nghĩa trong không gian đặc trưng hoặc không gian mô hình, sức mạnh của những bản vá thị giác thử nghiệm không được chú ý trong cơ chế chú ý vẫn bị bỏ qua. Trong bài báo này, chúng tôi học cách xác định các bản vá thị giác không liên quan đến ngữ nghĩa ở không gian đầu vào và bối cảnh hóa chúng về mặt ngữ nghĩa, do đó cho phép chúng truyền bá thông tin ngữ nghĩa đến những bản vá quan trọng đó trong suốt các giai đoạn học đặc trưng thị giác.
